\scanpage{35}

pour les $i<n$. \deleted{Prouvons que} Nous savons donc déjà
que $(\unsure{\alpha})$ est bijectif pour $i\leq n-1$,
injectif pour $i=n$, on veut gagner un cran.
On a une suite exacte
\diagram{\begin{tikzcd}[column sep=small]
{}\arrow[r,dash] & F\arrow[r] & L\arrow[r] & G\arrow[r] &0
\end{tikzcd}}
où \inserted{$G$ loc libre,} $L$ est une somme directe de
modules loc. libres \deleted{\unsure{sommes}} $\simeq L_\alpha$,
donc les assertions à prouver sont vraies pour $L$.
Considérons l'hom de suites exactes
\diagram{\begin{tikzcd}[column sep=small]
H^i(L)\arrow[r,dash]\arrow[d,"\alpha"',"\sim"] &
H^i(G)\arrow[r,dash]\arrow[d,"\beta"] &
H^{i+1}(F)\arrow[r]\arrow[d,"\varphi"] &
H^{i+1}(L)\arrow[r,dash]\arrow[d,"\alpha'"',"\sim"] &
H^{i+1}(G)\arrow[d,"\beta'"] \\
H^i(\hat L)\arrow[r] & H^i(\hat G)\arrow[r] &
H^{i+1}(\hat F)\arrow[r,dash] & H^{i+1}(\hat L)\arrow[r] &
H^{i+1}(\hat G)
\end{tikzcd}}
pour $i=n-1,n$ donc $i+1=n,n+1$.
Dans le premier cas, on veut prouver que la flèche médiane
est surjective, dans le deuxième qu'elle est injective.

a) $i=n-1$, on sait par hyp. de réc. que $\beta$ est
surjective \deleted{(ou bijective)}
\deleted{il en est de même de $\alpha'$},
\inserted{de plus $\alpha'$ est surj par hyp.}
\deleted{\unsure{et puis par l'hyp. de réc.}}
\inserted{$\beta'$ est injective.}
Par lemme des 5, $\varphi$ est \deleted{bijective} surjective
\inserted{(donc bijective)}.

b) $i=n$, On sait par \deleted{a) \unsure{que}
\unsure{l'hyp} \unsure{pour} $\beta$}
récurrence que $\alpha$ et $\beta$ sont bijectives,
donc $\alpha$ est surj, $\beta$ inj, de plus $\alpha'$ est inj.

\margin{34}
